% Processing steps of a simple web server for static content; the bars
% mark the steps that may wait for the network or for the disk.
% Usage: % Processing steps of a simple web server for static content; the bars
% mark the steps that may wait for the network or for the disk.
% Usage: % Processing steps of a simple web server for static content; the bars
% mark the steps that may wait for the network or for the disk.
% Usage: % Processing steps of a simple web server for static content; the bars
% mark the steps that may wait for the network or for the disk.
% Usage: \input{figures/l06-web-steps}   (width ~116 mm)
\begin{tikzpicture}[x=1mm, y=1mm, font=\scriptsize\sffamily,
  st/.style={draw, rounded corners=1.5pt, fill=ln-box, align=center,
    minimum width=14.6mm, minimum height=9mm, inner sep=0.5pt},
  num/.style={font=\tiny\sffamily, text=ln-muted, anchor=south},
  net/.style={draw=none, fill=ln-accent!60},
  dsk/.style={draw=none, fill=ln-caution!60},
  >={Stealth[length=1.2mm]}]
  \node[st] (s1) at (7.3,0)   {\iflangja{接続の\\受け付け}{accept\\connection}};
  \node[st] (s2) at (24.2,0)  {\iflangja{要求の\\読み込み}{read\\request}};
  \node[st] (s3) at (41.1,0)  {\iflangja{要求の\\解析}{parse\\request}};
  \node[st] (s4) at (58.0,0)  {\iflangja{ファイル\\の検索}{find\\file}};
  \node[st] (s5) at (74.9,0)  {\iflangja{ヘッダ\\の計算}{compute\\header}};
  \node[st] (s6) at (91.8,0)  {\iflangja{ヘッダ\\の送信}{send\\header}};
  \node[st] (s7) at (108.7,0) {\iflangja{ファイル\\読込・送信}{read file,\\send data}};
  \foreach \i [evaluate=\i as \j using int(\i+1)] in {1,...,6} {
    \draw[->] (s\i.east) -- (s\j.west);
  }
  \foreach \i in {1,...,7} { \node[num] at (s\i.north) {\i}; }
  % may wait for the network
  \foreach \x in {7.3,24.2,91.8,108.7} {
    \fill[net] (\x-7.3,-6.2) rectangle (\x+7.3,-5.0);
  }
  % may wait for the disk
  \foreach \x in {58.0,108.7} {
    \fill[dsk] (\x-7.3,-8.2) rectangle (\x+7.3,-7.0);
  }
  % legend
  \fill[net] (22,-12.6) rectangle (27,-11.4);
  \node[anchor=west] at (28,-12) {\iflangja{ネットワークを待つ可能性}{may wait for the network}};
  \fill[dsk] (68,-12.6) rectangle (73,-11.4);
  \node[anchor=west] at (74,-12) {\iflangja{ディスクを待つ可能性}{may wait for the disk}};
\end{tikzpicture}
   (width ~116 mm)
\begin{tikzpicture}[x=1mm, y=1mm, font=\scriptsize\sffamily,
  st/.style={draw, rounded corners=1.5pt, fill=ln-box, align=center,
    minimum width=14.6mm, minimum height=9mm, inner sep=0.5pt},
  num/.style={font=\tiny\sffamily, text=ln-muted, anchor=south},
  net/.style={draw=none, fill=ln-accent!60},
  dsk/.style={draw=none, fill=ln-caution!60},
  >={Stealth[length=1.2mm]}]
  \node[st] (s1) at (7.3,0)   {\iflangja{接続の\\受け付け}{accept\\connection}};
  \node[st] (s2) at (24.2,0)  {\iflangja{要求の\\読み込み}{read\\request}};
  \node[st] (s3) at (41.1,0)  {\iflangja{要求の\\解析}{parse\\request}};
  \node[st] (s4) at (58.0,0)  {\iflangja{ファイル\\の検索}{find\\file}};
  \node[st] (s5) at (74.9,0)  {\iflangja{ヘッダ\\の計算}{compute\\header}};
  \node[st] (s6) at (91.8,0)  {\iflangja{ヘッダ\\の送信}{send\\header}};
  \node[st] (s7) at (108.7,0) {\iflangja{ファイル\\読込・送信}{read file,\\send data}};
  \foreach \i [evaluate=\i as \j using int(\i+1)] in {1,...,6} {
    \draw[->] (s\i.east) -- (s\j.west);
  }
  \foreach \i in {1,...,7} { \node[num] at (s\i.north) {\i}; }
  % may wait for the network
  \foreach \x in {7.3,24.2,91.8,108.7} {
    \fill[net] (\x-7.3,-6.2) rectangle (\x+7.3,-5.0);
  }
  % may wait for the disk
  \foreach \x in {58.0,108.7} {
    \fill[dsk] (\x-7.3,-8.2) rectangle (\x+7.3,-7.0);
  }
  % legend
  \fill[net] (22,-12.6) rectangle (27,-11.4);
  \node[anchor=west] at (28,-12) {\iflangja{ネットワークを待つ可能性}{may wait for the network}};
  \fill[dsk] (68,-12.6) rectangle (73,-11.4);
  \node[anchor=west] at (74,-12) {\iflangja{ディスクを待つ可能性}{may wait for the disk}};
\end{tikzpicture}
   (width ~116 mm)
\begin{tikzpicture}[x=1mm, y=1mm, font=\scriptsize\sffamily,
  st/.style={draw, rounded corners=1.5pt, fill=ln-box, align=center,
    minimum width=14.6mm, minimum height=9mm, inner sep=0.5pt},
  num/.style={font=\tiny\sffamily, text=ln-muted, anchor=south},
  net/.style={draw=none, fill=ln-accent!60},
  dsk/.style={draw=none, fill=ln-caution!60},
  >={Stealth[length=1.2mm]}]
  \node[st] (s1) at (7.3,0)   {\iflangja{接続の\\受け付け}{accept\\connection}};
  \node[st] (s2) at (24.2,0)  {\iflangja{要求の\\読み込み}{read\\request}};
  \node[st] (s3) at (41.1,0)  {\iflangja{要求の\\解析}{parse\\request}};
  \node[st] (s4) at (58.0,0)  {\iflangja{ファイル\\の検索}{find\\file}};
  \node[st] (s5) at (74.9,0)  {\iflangja{ヘッダ\\の計算}{compute\\header}};
  \node[st] (s6) at (91.8,0)  {\iflangja{ヘッダ\\の送信}{send\\header}};
  \node[st] (s7) at (108.7,0) {\iflangja{ファイル\\読込・送信}{read file,\\send data}};
  \foreach \i [evaluate=\i as \j using int(\i+1)] in {1,...,6} {
    \draw[->] (s\i.east) -- (s\j.west);
  }
  \foreach \i in {1,...,7} { \node[num] at (s\i.north) {\i}; }
  % may wait for the network
  \foreach \x in {7.3,24.2,91.8,108.7} {
    \fill[net] (\x-7.3,-6.2) rectangle (\x+7.3,-5.0);
  }
  % may wait for the disk
  \foreach \x in {58.0,108.7} {
    \fill[dsk] (\x-7.3,-8.2) rectangle (\x+7.3,-7.0);
  }
  % legend
  \fill[net] (22,-12.6) rectangle (27,-11.4);
  \node[anchor=west] at (28,-12) {\iflangja{ネットワークを待つ可能性}{may wait for the network}};
  \fill[dsk] (68,-12.6) rectangle (73,-11.4);
  \node[anchor=west] at (74,-12) {\iflangja{ディスクを待つ可能性}{may wait for the disk}};
\end{tikzpicture}
   (width ~116 mm)
\begin{tikzpicture}[x=1mm, y=1mm, font=\scriptsize\sffamily,
  st/.style={draw, rounded corners=1.5pt, fill=ln-box, align=center,
    minimum width=14.6mm, minimum height=9mm, inner sep=0.5pt},
  num/.style={font=\tiny\sffamily, text=ln-muted, anchor=south},
  net/.style={draw=none, fill=ln-accent!60},
  dsk/.style={draw=none, fill=ln-caution!60},
  >={Stealth[length=1.2mm]}]
  \node[st] (s1) at (7.3,0)   {\iflangja{接続の\\受け付け}{accept\\connection}};
  \node[st] (s2) at (24.2,0)  {\iflangja{要求の\\読み込み}{read\\request}};
  \node[st] (s3) at (41.1,0)  {\iflangja{要求の\\解析}{parse\\request}};
  \node[st] (s4) at (58.0,0)  {\iflangja{ファイル\\の検索}{find\\file}};
  \node[st] (s5) at (74.9,0)  {\iflangja{ヘッダ\\の計算}{compute\\header}};
  \node[st] (s6) at (91.8,0)  {\iflangja{ヘッダ\\の送信}{send\\header}};
  \node[st] (s7) at (108.7,0) {\iflangja{ファイル\\読込・送信}{read file,\\send data}};
  \foreach \i [evaluate=\i as \j using int(\i+1)] in {1,...,6} {
    \draw[->] (s\i.east) -- (s\j.west);
  }
  \foreach \i in {1,...,7} { \node[num] at (s\i.north) {\i}; }
  % may wait for the network
  \foreach \x in {7.3,24.2,91.8,108.7} {
    \fill[net] (\x-7.3,-6.2) rectangle (\x+7.3,-5.0);
  }
  % may wait for the disk
  \foreach \x in {58.0,108.7} {
    \fill[dsk] (\x-7.3,-8.2) rectangle (\x+7.3,-7.0);
  }
  % legend
  \fill[net] (22,-12.6) rectangle (27,-11.4);
  \node[anchor=west] at (28,-12) {\iflangja{ネットワークを待つ可能性}{may wait for the network}};
  \fill[dsk] (68,-12.6) rectangle (73,-11.4);
  \node[anchor=west] at (74,-12) {\iflangja{ディスクを待つ可能性}{may wait for the disk}};
\end{tikzpicture}
